% Transcription — fonds Alexandre Grothendieck, shelfmark 75, batch 2 (pages 21-32).
% Established from the facsimile archives/batches/75/batch-02.pdf (prepared
% from a copy of 75.pdf uploaded by the user; PDF page k = archivists' page
% 20+k, no cover sheet), per the skill transcribe-grothendieck. The folder is
% thirty-two pages; this is the second and last batch. Batch 1 (pages 1-20)
% is transcribed separately and did not exist when this pass was made, so no
% notation could be carried over from it.
%
% The dating is the folder's own, copied verbatim; the group « Géométrie et
% topologie combinatoire » (69-89) holds it. Folder 75 is « Cours C4, DEA :
% notes manuscrites (s.d.), lettre (1978), tapuscrit (1977-1978) ».
%
% Pages skipped: 28 and 29 — two copies (the first a faint carbon, the second
% a photocopy with a pencil scrawl across it) of a typed letter signed by
% Grothendieck, Villecun 30.01.78, to the Minister of Education, on the
% ordonnance of 2.11.1945 — the « lettre (1978) » of the inventory. It is his,
% but it carries no mathematics, so by the skill it is skipped and not
% described in the body.
%
% Pages summarised: none. Pages 30-32 are a typescript, but it is his own and
% not a published text: two announcements of his 1977-78 courses, headed
% « COURS C4 1977-78 de A. GROTHENDIECK » (pp. 30-31, « Géométrie de
% l'icosaèdre ») and « Cours DEA 1977/78 de A. Grothendieck » (p. 32,
% « Géométrie arithmétique des polygônes réguliers »). As his own typescript
% they are transcribed in full. The printed letterhead of the Institut de
% Mathématiques on p. 30 (with « Montpellier, le » left blank) is paper, not
% text, and is omitted. On pp. 30-31 the gothic A_5 and S_5 were inked by hand
% into blanks; the double-struck F_4, F_5 look typed. On p. 32 a pencil « ! »
% in the left margin and a pencil underline of « en commun » are recorded in a
% note without attribution — nothing in the stroke says whose they are.
%
% Pages transcribed from his hand: 21-27, in black ink, fairly fast but
% mostly legible; the formulas carry throughout, a few connecting words do not.
% His pagination: page 21 opens with a circled « 8 » — a section number of his,
% presumably continuing a run begun in batch 1.
%
% What the batch is about. Preparatory notes for the courses, around torsors
% and orientations. P. 21 (« 8 », « Digression sur structures cubiques »): over
% a ring A with 2 ≠ 0, the category of free modules E with a direct-sum
% decomposition into n rank-one pieces E_i, each with a « bi-base » {e_i, f_i},
% e_i + f_i = 0 — equivalently a subset B of 2n elements stable under x -> -x
% whose representatives mod ±1 form a basis — is equivalent to the category of
% sets B of cardinal 2n with a fixed-point-free involution σ; the inverse
% functor takes E = ⊕ E_i with E_i = B_i ∧^{±1} (A with ±1 acting by
% homotheties). Pp. 22-23: a (combinatorial) contour C of length n determines
% its two-element set of orientations ω and a torsor under G+ = Aut+(C) ≃
% ω ∧_{±1} Z_n; conversely a pair (ω, T) determines a contour; the
% « Exer (DEA) » generalises: for a semidirect product G = Z.Γ, a G-torsor
% amounts to a Γ-torsor ω and a torsor T_0 under the twisted group ω ∧_Γ Z.
% Pp. 24-25 (« Orientations »): a) real vector space E, ω_E = P ×_G {±1, s} ≃
% P/Gl+(n,R) with P = Isom(R^n, E); b) free quadratic modules over k with 2
% invertible and an orthonormal basis, ω_E = P ×_G (±1, s) ≃ P/SO(n,k);
% compatibility by transitivity of extension of the structure group along
% O(n,k) -> Gl(n,k) -> {±1}; c) a finite set I, ω_I = P ×_{S_n} {±1} ≃ P/A_n
% with P = numérot(I); compatibility with b) through k^I with its standard
% scalar product and S_n -> O(n,k) -> {±1}, hence ω_{E(I,R)} ≃ ω_I. Pp. 26-27
% (« Formules d'associativité »): for E = ⊕ E_i real, ω_E ≃ ω_{I^-} ∧ ⋀ ω_{E_i}
% with I^- the indices of odd dimension (boxed), with the verification by
% ordered bases; the same for quadratic modules; for a map J -> I of finite
% sets, ω_J ≃ ω_{I^-} ∧ ⋀ ω_{J_i}, I^- the odd fibres; « Compatibilité entre
% les isom. can. a) b) c) »; « Cas d'extensions »: for a flag 0 = E_0 ⊂ … ⊂
% E_d = E, define ω_E ≃ ⋀ ω_{E_i/E_{i-1}}. Pp. 30-32: the two course
% announcements (icosahedron combinatorially and over arbitrary fields, the
% dictionary with the oriented five-element set, A_5 ≃ SL(2,F_4) ≃ …, S_5 ≃
% GP(1,F_5) ≃ SO(3,F_5); regular polygons over fields and rings, the
% « general » icosahedron over a ring finite over Z).
%
% Notation fixed once for the batch: 𝒞, 𝒞' as \mathcal{C}, \mathcal{C}' on p. 21;
% the contracted product written with ∧ and the group as subscript, as he
% writes it (\omega \wedge_{\{\pm 1\}} \mathbf{Z}_n); the fibre product as
% P \times_G …; his Gl, GP, SO, O kept as written; the symmetric and alternating
% groups \mathfrak{S}_n, \mathfrak{A}_n; his double-struck Z, R, N, F as
% \mathbf{Z}, \mathbf{R}, \mathbf{N}, \mathbb{F}; the Z of p. 23 is a plain Z,
% as written. His « = » for « même » on p. 27 is kept as the sign.
%
% Readings to check first: « Digression » in the heading of p. 21; « sous- »
% before « modules » and « donc tt » on p. 21; the word before
% « orientations » and « (sinon compris) » on p. 22; the two overwritten
% letters « Z et Γ » on p. 23; « ord. » in the margin of p. 25 and the whole
% last line of p. 25; « quotient » and the insertion « donc définit le même
% él. » on p. 26; « (surjection si on veut) » and the blotted J_i on p. 27.
%
% Pass: Opus 5.5 (claude-opus-5-5), 2026-09-24 — first pass, unchecked against the pages by a human.

\documentclass[11pt,a4paper]{article}
% Le préambule commun à toutes les transcriptions du fonds Grothendieck.
%
% Il tient en une idée : **l'incertitude se compose, elle ne se résout pas.**
% Une transcription qui lisse un mot douteux a détruit la seule chose pour
% laquelle on la faisait — un lecteur doit pouvoir savoir, à chaque endroit,
% ce qui était sur le feuillet et ce qui a été ajouté. Les six macros qui
% suivent sont donc l'appareil critique, et non de la décoration.
%
% Les mêmes commandes sont reconnues par `scripts/render.mjs`, qui produit la
% vue de lecture du site. Une macro ajoutée ici doit l'être aussi là-bas,
% faute de quoi le rendu échouera bruyamment — ce qui est voulu : un rendu
% silencieusement faux est pire qu'un rendu absent.

\ProvidesPackage{grothendieck}[2026/08/08 Transcriptions du fonds Grothendieck]

\usepackage{amsmath,amssymb,amsthm}
\usepackage{mathrsfs}
\usepackage{stmaryrd}
% Les diagrammes commutatifs sont partout dans ce fonds : c'est la langue
% même de la géométrie algébrique de Grothendieck, et une transcription qui
% les rendrait en prose ne transcrirait rien.
\usepackage{tikz-cd}
% Français par défaut — c'est la langue des feuillets — mais l'anglais est
% chargé aussi : la traduction et le résumé sont des documents anglais, et la
% typographie française (espaces avant les deux-points) y serait une faute.
% Ces documents basculent par \selectlanguage{english} en tête de corps.
\usepackage[english,french]{babel}
\usepackage{fontspec}
\usepackage{geometry}
\geometry{a4paper,margin=25mm,marginparwidth=28mm}
\usepackage{marginnote}
\usepackage{xcolor}
% ulem, not soul. `soul`'s \st and \ul are built on a character-by-character
% scan that XeLaTeX's font machinery breaks: the document fails with an
% undefined control sequence rather than degrading. ulem does the same two jobs
% and is Unicode-engine safe. `normalem` stops it hijacking \emph, which the
% transcriptions use for Grothendieck's own emphasis.
\usepackage[normalem]{ulem}
\usepackage{hyperref}
\hypersetup{colorlinks=true,linkcolor=black,urlcolor=black,pdfborder={0 0 0}}

% --- Métadonnées du lot -----------------------------------------------------
% Reprises telles quelles par le script de rendu, qui les affiche en tête de
% la vue de lecture. Elles fixent ce que le fichier prétend couvrir : sans
% elles, une transcription flotte sans point d'attache dans un fonds de seize
% mille feuillets.

\newcommand{\folder}[1]{\gdef\grothFolder{#1}}
\newcommand{\batch}[1]{\gdef\grothBatch{#1}}
\newcommand{\leaves}[2]{\gdef\grothFirst{#1}\gdef\grothLast{#2}}
\newcommand{\foldertitle}[1]{\gdef\grothFoldertitle{#1}}
\gdef\grothFolder{?}\gdef\grothBatch{?}\gdef\grothFirst{?}\gdef\grothLast{?}\gdef\grothFoldertitle{}

% --- Le résumé --------------------------------------------------------------
%
% Il ouvre la lecture modernisée, sous le titre « Résumé », et il est composé
% en petit. Ce n'est pas un ornement : c'est le seul passage du document qui
% ne soit pas des mathématiques, et un lecteur doit pouvoir le reconnaître
% avant de l'avoir lu — puis le sauter s'il connaît déjà le sujet, ou s'y
% arrêter s'il ne le connaît pas. Le filet qui le suit marque la reprise du
% corps.

\newenvironment{resume}
  {\small\setlength{\parskip}{0.45em}}
  {\par\medskip\hrule\medskip}

% --- Mots-clés ---------------------------------------------------------------
%
% En anglais, en clôture du résumé de la lecture modernisée : le vocabulaire
% moderne sous lequel chercher ce que les feuillets construisent. Le manifeste
% du site (`npm run manifest`) les extrait pour étiqueter la cote entière —
% cette ligne est la source unique des tags, il n'y a pas de fichier à côté.
\newcommand{\keywords}[1]{\par\smallskip\noindent
  {\footnotesize\textsc{Keywords} --- \textit{#1}}\par}

% --- L'appareil critique ----------------------------------------------------

% Le numéro de feuillet, en marge. C'est l'ancre qui permet de régler un
% désaccord sur une formule en regardant *un* feuillet plutôt qu'une chemise
% de six cents. Le site s'en sert aussi pour tourner les pages du fac-similé
% quand on fait défiler la transcription.
\newcommand{\leaf}[1]{%
  \marginnote{\footnotesize\color{gray!70}\textsf{#1}}%
  \label{leaf:#1}\ignorespaces}

% Illisible : on ne devine pas. Un mot inventé qui se lit comme les autres est
% le pire résultat possible.
\newcommand{\ill}{\textcolor{red!60!black}{[\,\ldots\,]}}

% Lecture douteuse : le mot est donné, et signalé comme incertain.
\newcommand{\uncertain}[1]{\uline{#1}}

% Ajout de l'éditeur — une lettre manquante, un mot sous-entendu.
\newcommand{\add}[1]{\textcolor{blue!50!black}{[#1]}}

% Biffé par Grothendieck lui-même. Ce qu'il a rayé fait partie du document :
% une rature dit souvent où la pensée a changé de direction.
\newcommand{\struck}[1]{\sout{#1}}

% Note du transcripteur, sur l'état matériel du feuillet.
\newcommand{\note}[1]{\par\smallskip{\small\color{gray!60!black}\textit{[#1]}}\par\smallskip}

% Note marginale de Grothendieck, distincte du corps du texte.
\newcommand{\marginal}[1]{\par\smallskip\noindent\hspace{1em}%
  {\small\color{gray!60!black}#1}\par\smallskip}

% --- Titre ------------------------------------------------------------------

\AtBeginDocument{%
  \begin{center}
    {\large\bfseries Fonds Alexandre Grothendieck --- cote n\textsuperscript{o}~\grothFolder}\\[2pt]
    {\itshape\grothFoldertitle}\\[4pt]
    {\small feuillets \grothFirst--\grothLast\ (lot \grothBatch)}\\[2pt]
    {\footnotesize\color{gray!60!black}Transcription établie d'après le fac-similé numérisé
     par l'Université de Montpellier.\\
     Les lectures incertaines sont soulignées, les illisibles notées [\,\ldots\,],
     les ajouts entre crochets.}
  \end{center}
  \vspace{1em}}

\endinput


\folder{75}
\batch{2}
\pages{21}{32}
\dating{1977-1978}
\watermark{Édition de démonstration}
\foldertitle{Cours C4, DEA : notes manuscrites (s.d.), lettre (1978), tapuscrit (1977-1978)}

\begin{document}

\section*{\uncertain{Digression} sur structures cubiques}

\page{21}
\note{le titre est précédé, dans l'original, d'un « 8 » cerclé : numéro de
section de l'auteur}

$A$ anneau avec $2 \neq 0$, $n \in \mathbf{N}^*$.

$\mathcal{C}$ cat.\ modules
libres $E$ sur $A$, avec \struck{famille} ens.\ de $n$ \uncertain{sous-}modules
libres $E_i$ ($i \in I$, $\operatorname{card} I = n$) de somme directe $E$, et
chaque $E_i$ muni d'une bi-base $\{e_i, f_i\}$ ($e_i + f_i = 0$)
\note{« cat. » et « $n$ » sont ajoutés au-dessus de la ligne}

[ou encore, $E$ muni d'une partie $B$ de card.\ $2n$ stable par $x \mapsto -x$,
et telle que un -- \uncertain{donc tt} -- syst.\ de représentants de $B/\pm 1$
dans $B$ soit une base]
\note{« de card. $2n$ » est ajouté au-dessus de la ligne}

\[ \wr\!\downarrow \]

$\mathcal{C}'$ cat.\ des ens.\ $B$ de cardinal $2n$ munis d'une rel.\
d'équivalence \struck{\ill{}} à orbites de card.\ $2$, ou encore muni d'une
permutation $\sigma$ involution ($\sigma^2 = \mathrm{id}$) sans pt fixe.
\note{« $B$ » est ajouté au-dessus de la ligne}

$\mathcal{C} \to \mathcal{C}'$ évident

$\mathcal{C}' \to \mathcal{C}$ : si $I = B/\sigma$, on considère
\[
  E = \bigoplus_{i \in I} E_i,
\]
où
\[
  E_i = B_i \wedge^{\{\pm 1\}} \ldots
\]
\note{après le signe $\wedge$, un symbole est surchargé et barré ; il n'est
pas lu}
où \uncertain{$D$} est le $A$-module \uncertain{$A_g$} sur lequel $\pm 1$ opère
par homothéties.

\section*{Contours}

\page{22}
Contour (combin.) convexe $C$ définit
\note{« (combin.) » est ajouté au-dessus de la ligne ; le mot qui suit,
lu « convexe », paraît remplacé par lui sans être nettement barré}

\begin{enumerate}
  \item L'ens.\ $\omega = \omega(C)$ à deux éléments de ses deux orientations.
\end{enumerate}

À partir de là, on peut déjà écrire \uncertain{(sinon compris)} le groupe des
translations du contour, par
\[
  G^+ = \operatorname{Aut}^+(C) \simeq \omega \wedge_{\{\pm 1\}} \mathbf{Z}_n
\]
(où $\{\pm 1\}$ opère sur $\mathbf{Z}_n$ par homothéties)

\begin{enumerate}
  \setcounter{enumi}{1}
  \item Un torseur sous $G^+$.
\end{enumerate}

\struck{Exercice} Et inversement, les couples formés par $(\omega, T)$, où
$\omega$ un \struck{\ill{}} ens.\ à $2$ él.\ et $T$ un torseur sous
$G^+ = \omega \wedge_{\{\pm 1\}} \mathbf{Z}_n$, définit un

\page{23}
contour.

\textit{Exer} (DEA) Soit $G = Z.\Gamma$ produit ½-direct de \uncertain{$Z$} et
\uncertain{$\Gamma$} ($\Gamma$ opère sur $Z$). Montrer que la donnée d'un
torseur $T$ sous $G$ équivaut à celle de
\note{les deux lettres après « ½-direct de » sont surchargées, l'une sur
l'autre ; l'ordre n'est pas sûr. « $T$ » est ajouté au-dessus de « torseur »}

\begin{itemize}
  \item[a)] d'un torseur $\omega$ sous $\Gamma$
  \item[b)] d'un torseur $T_0$ sous ${}^{\omega}Z = \omega \wedge_{\Gamma} Z$
\end{itemize}

\section*{Orientations}

\page{24}
a) Vectoriel réel $E$ de dim.\ finie $n$. La donnée de $E$ équivaut à celle du
torseur à dr.
\[
  P = \operatorname{Isom}(\mathbf{R}^n, E) \simeq \text{bases}(E)
\]
sous
\[
  G = \mathrm{Gl}(n, \mathbf{R}).
\]
On a un hom.

\begin{tikzcd}[column sep=small]
  G = \mathrm{Gl}(n, \mathbf{R}) \arrow[rr, "s"] \arrow[dr, "\det"'] & & \{\pm 1\} \\
  & \mathbf{R}^* \arrow[ur, "\operatorname{sgn}"'] &
\end{tikzcd}

$u \mapsto$ \struck{$\det$} $\operatorname{sgn} \det u$, noyau
$\mathrm{Gl}^+(n, \mathbf{R})$, et
\[
  \omega_E = P \times_G \{\pm 1, s\} \simeq P / \mathrm{Gl}^+(n, \mathbf{R})
\]
est l'ens.\ des orientations de $E$.

b) Modules quadratiques libres \struck{\ill{}} $E$ sur un anneau comm.\ $k$ où
$2$ est inversible [$\exists$ base orthonormale, \struck{\ill{}} $e_i^2 = 1$,
$e_i e_j = 0$ si $i \neq j$]. La donnée de $E$ de rang $n$ équivaut à celle du
tors.\ à droite
\note{« quadratiques » est ajouté au-dessus de la ligne ; « $E$ » remplace un
mot barré}
\[
  P = \operatorname{Isom}_{k\text{-quadr.}}(k^n, E) \simeq \text{bases orthon.}(E)
\]
sous
\[
  G = \mathrm{O}(n, k) = \operatorname{Aut}_{k\text{-quadr.}}(k^n, Q_n)
  \subset \mathrm{Gl}(n, k)
\]
Or ici
\[
  G = \mathrm{O}(n, k) \xrightarrow{\ \det = s\ } \{\pm 1\}
\]
(noyau $\mathrm{SO}(n, k)$)
et l'ens.\ des orientations de $E$ est par définition
\[
  \omega_E = P \times_G (\pm 1, s) \simeq P / \mathrm{SO}(n, k).
\]

\textit{Compatibilité entre a) et b)} Résulte de la transitivité de l'opération
d'ext.\ du groupe d'opérateurs dans un torseur, et du diagramme comm.

\begin{tikzcd}
  \mathrm{O}(n, k) \arrow[r, hook, "\mathrm{incl}"] \arrow[rr, bend right=20, "\det"'] &
  \mathrm{Gl}(n, k) \arrow[r, "\operatorname{sgn} \det"] & \{\pm 1\}
\end{tikzcd}

\page{25}
c) Orientation d'un ensemble fini $I$ (de card.\ $n$). La donnée de $I$
équivaut à celle d'un torseur
\note{« (de card. $n$) » est ajouté au-dessus de la ligne}
\[
  P = \operatorname{Isom}(I_n, I) = \text{numérot}(I)
\]
\note{avant $I_n$, une première écriture de l'intervalle $\{1, \ldots, n\}$ est
barrée}
($I_n = \{ m \in \mathbf{N}^* \mid m \leq n \} = \{1, 2, \ldots, n\}$) sous le
groupe
\[
  \mathfrak{S}_n = \mathfrak{S}_{I_n} = \operatorname{Aut}(I_n).
\]
\note{une première lettre, raturée, précède $\mathfrak{S}_n$}
On a un hom.\ canonique « signature »
\[
  \mathfrak{S}_n \xrightarrow{\ s = \operatorname{sgn}\ } \{\pm 1\}
\]
noyau $\mathfrak{A}_n$
et par définition l'ens.\ des $2$ orientations de $I$ est
\[
  \omega_I = P \times_{\mathfrak{S}_n} \{\pm 1\} \simeq P / \mathfrak{A}_n
\]
\note{un premier membre, raturé, précède $P \times_{\mathfrak{S}_n}$}

\textit{Compatibilités avec a) et b)}

\textit{Avec} b) On a sur $E_{(I,k)} = k^I$ un produit scalaire canonique
\[
  (x, y) = \sum_{i \in I} x_i y_i
\]
pour lequel $(e_i)$ est une base orthonormale. Ainsi $I$ \struck{\ill{}}
\uncertain{donne} naissance à un module quadratique $E(I, k)$. Ceci posé, on
définit une bijection canonique
\[
  \omega_{E(I,k)} \simeq \omega_I
\]
\note{la seconde lettre de $E(I,k)$ ressemble ici à un X ; le contexte
($E_{(I,k)} = k^I$) donne $k$}
\marginal{Si $i_1, \ldots, i_n$ est une numérotation de $I$,
$e_{i_1}, \ldots, e_{i_n}$ est une base \uncertain{ord.}\ orthon.\ de
$E$, \ldots}
par transitivité dans

\begin{tikzcd}
  \mathfrak{S}_n \arrow[r, hook, "\mathrm{incl.}"] \arrow[rr, bend right=20, "\operatorname{sgn}"'] &
  \mathrm{O}(n, k) \arrow[r, "\det"] & \{\pm 1\}
\end{tikzcd}

En particulier, on trouve (comp.\ avec a)),
\[
  \omega_{E(I, \mathbf{R})} \simeq \omega_I
\]
$\to$ \uncertain{considéré} \ill{} vect.\ réel \uncertain{se pense}
\note{ligne ajoutée en bas du feuillet, sous et à droite de la formule ; lecture
très incertaine}

\section*{Formules d'associativité}

\page{26}
a) Vectoriels réels. \struck{Soit} $I$ un ens.\ fini, $(E_i)_{i \in I}$ famille
de vectoriels réels de dim finie, $E = \coprod_{i \in I} E_i$ la somme directe
($=$ produit direct). \struck{Soit} $I^+$ ($I^-$) l'ens.\ des $i \in I$ tels
que $d_i = \dim_{\mathbf{R}} E_i$ soit pair (impair). On a un isom.\ canonique
\note{« réels » est ajouté au-dessus de la ligne ; le $\dim$ surcharge une
première écriture}
\[
  \boxed{\ \omega_E \simeq \omega_{I^-} \wedge \bigwedge_{i \in I} \omega_{E_i}\ }
\]
obtenu ainsi, en regardant le $2^{\text{ième}}$ mb comme \uncertain{quotient}
de
\[
  \text{numérot}(I^-) \times \text{numérot}(I^+) \times
  \prod_{i \in I} \text{bases}(E_i)
\]
\note{« $\times\, \text{numérot}(I^+)$ » est ajouté au-dessus de la ligne et
rattaché par un trait après $\text{numérot}(I^-)$}
\marginal{NB Un ordre total de $I$ en induit un de $I^-$, donc
$\text{numérot}(I) \to \text{numérot}(I^-)$}
Soit $\{i_1, \ldots, i_{n^-}\}$ une num.\ de $I^-$, $j_1 \ldots j_{n^+}$ une
num.\ de $I^+$, \struck{et} $\forall i \in I$,
$B_i = (e_{i,1}, \ldots, e_{i,d_i})$ une base ordonnée de $E_i$, alors
\note{« $j_1 \ldots j_{n^+}$ une num. de $I^+$, » est ajouté au-dessus de la
ligne}
\[
  e_{i_1,1}, \ldots, e_{i_1,d_1},\ e_{i_2,1}, \ldots, e_{i_2,d_2},\ \ldots,\
  e_{i_n,1}, \ldots, e_{i_n,d_{i_n}}
\]
est une base de $E$. Si on change une $B_i$ par une matrice de
$\mathrm{Gl}(d_i, \mathbf{R})^+$, alors la base se trouve changée par une
matrice de $\mathrm{Gl}(d, \mathbf{R})^+$ ($d = \sum d_i = \operatorname{rg} E$),
\uncertain{donc définit le même} él.\ de $\omega_E$ ; si on change l'ordre de
$I^-$ par une permutation $\sigma^-$ (une transposition entre $\alpha$ et
$\alpha + 1$, disons \ldots) \uncertain{on change} la base par une matrice de
dét \struck{\ill{}} $\operatorname{sgn} \sigma^-$ ; par contre, si on change
\struck{par} l'ordre de $I^+$ par un $\sigma^+$, la matrice de passage est de
dét $+1$ en tous cas \ldots
\note{« donc définit le même él. de $\omega_E$ » est ajouté au-dessus de la
ligne}

\page{27}
b) Contexte des familles finies de modules quadratiques : $=$ formule, avec $=$
démonstration.

c) $J \xrightarrow{\ p\ } I$ application \uncertain{(surjection si on veut)}
d'ens.\ finis,
\[
  I^- = \{ i \in I \mid \operatorname{card} p^{-1}(i) \ \mathrm{impair} \}.
\]
\struck{\ill{}} [p.\ ex.\ $J$ un ens.\ muni d'une relation d'équivalence
$\mathcal{R}$, et $I = J/\mathcal{R}$]. Alors
\[
  \omega_J \simeq \omega_{I^-} \wedge \bigwedge_{i \in I} \omega_{J_i}
\]
\note{l'indice du dernier $\omega$ est empâté par une tache d'encre ; $J_i$ est
une lecture}

\textit{Compatibilité entre les isom.\ can.\ a) b) c)}

\section*{Cas d'extensions}

\[
  (0) = E_0 \subset E_1 \subset E_2 \cdots \subset E_d = E
\]
\struck{\ill{}} Définir un isom.\ canonique
\[
  \omega_E \simeq \bigwedge_{1 \leq i \leq d} \omega_{E_i / E_{i-1}}
\]

\section*{COURS C4 1977-78 de A.~GROTHENDIECK}

\page{30}
\note{tapuscrit de l'auteur, sur papier à en-tête de l'Institut de
Mathématiques de l'USTL, dont la date est restée en blanc ; les lettres
gothiques sont ajoutées à la main dans des blancs de la frappe}

\textit{Géométrie de l'icosaèdre}

Il est étrange que l'icosaèdre, étudié depuis plus de deux mille ans dans des
contextes et des optiques très divers, ne semble jamais avoir été étudié
systématiquement du point de vue combinatoire, ni de celui des réalisations
géométriques sur des corps (voire des anneaux) généraux. Ce "cours" se voudrait
une investigation commune sur ce thème, en partant de l'intuition combinatoire
la plus naïve, exemplifiée sur des modèles en carton (avec un coloriage
convenable des arètes).

Nous aborderons donc la théorie de l'icosaèdre d'un point de vue combinatoire
d'abord, en explicitant dans le détail le dictionnaire entre la géométrie
combinatoire de l'icosaèdre orienté, et celle de l'ensemble "orienté" à cinq
éléments (les cinq couleurs coloriant les arètes !) -- ce qui donnera en même
temps une motivation et une inspiration géométrique pour étudier le groupe des
automorphismes commun aux deux situations, savoir le groupe alterné
$\mathfrak{A}_5$ (qui est aussi le plus petit groupe simple non abélien). La
question de construire des "réalisations géométriques" de l'icosaèdre
combinatoire, comme "polyèdre régulier" sur le corps des réels, ou sur tout
autre corps, sera l'occasion à la fois d'approfondir la compréhension de la
structure combinatoire, et de nous familiariser avec les notions de géométrie
vectorielle et affine (voire de géométrie projective) sur des corps de
caractéristique quelconque.

\page{31}
En relation avec l'icosaèdre, nous découvrirons que les corps de
caractéristique deux et cinq jouent un rôle très particulier, tant pour l'étude
des réalisations géométriques, que pour l'étude combinatoire abstraite, qui
peut s'expliciter par des notions de géométrie sur les corps $\mathbb{F}_4$ et
$\mathbb{F}_5$ -- ces deux dictionnaires étant subordonnés aux "isomorphismes
exceptionnels" :
\begin{align*}
  \mathfrak{A}_5 &\simeq \mathrm{Sl}(2, \mathbb{F}_4) \quad
    (\simeq \mathrm{GP}(1, \mathbb{F}_4) \simeq \mathrm{SO}(3, \mathbb{F}_4)) \\
  \mathfrak{S}_5 &\simeq \mathrm{GP}(1, \mathbb{F}_5) \quad
    (\simeq \mathrm{SO}(3, \mathbb{F}_5)) .
\end{align*}

\section*{Cours DEA 1977/78 de A.~Grothendieck}

\page{32}
\textit{Géométrie arithmétique des polygônes réguliers}

\note{au crayon, un « ! » dans la marge gauche en face de la première ligne,
et un trait sous « en commun » ; rien n'indique de quelle main}

Ce "cours" se fera dans l'esprit d'une investigation en commun de multiples
facettes de la notion de polygône régulier au-dessus d'un corps de base, puis
d'un anneau de base quelconque. Ce sera une occasion stimulante pour nous
familiariser avec les notions de base de géométrie vectorielle, affine,
projective -- avec ou sans métrique -- en dimensions un, deux et trois, et avec
les phénomènes particuliers à ces dimensions, ainsi qu'avec les aspects parfois
imprévus que prennent des notions familières (comme celle de réflexion, ou de
centre d'un polygône régulier) sur les corps de caractéristique positive, et
notamment sur les corps finis. Le besoin d'une géométrie sur des anneaux de
base généraux (tel l'anneau "sur lequel s'écrit" le n-polygône régulier "le
plus général") se fera sentir en cours de route, et nous amènera sans doute à
affleurer quelques points de fondements de géométrie algébrique. Si le temps le
permet et si d'autres investigations ne nous retiennent , nous tenterons, dans
la volée, de construire l'icosaèdre régulier "général" sur un anneau de base
convenable, fini sur l'anneau $\mathbf{Z}$ des entiers, ayant des corps
résiduels de toutes caractéristiques. (La caractéristique deux donnera sans
doute du fil à retordre !).

A partir d'un point de départ intuitif et familier, ce "cours" voudrait être
une ouverture à un univers fascinant, pleins de mystères qui ne demandent qu'à
être élucidés. Il ne peut donner des fruits qu'avec une participation active
de l'étudiant et un travail personnel important, notamment pour développer avec
le détail nécessaire les nombreux points (de fondements notamment) qui
n'auront été qu'esquissés dans le "cours". Il lui est très fortement conseillé
de suivre en même temps le cours C4 sur l'icosaèdre, qui se fera dans un esprit
analogue. Il y aura sans doute de nombreuses interférences entre ces deux cours
de réflexion.

\end{document}
